\FigureLayoutDeclare{img/2014/shaanxi_liberal/q15_fig1.png}{5.4cm}{15bp 17bp 14bp 7bp}

\chapter{2014年陕西卷（文）}

\section{单选题}

\begin{problem}
已知集合 \(M = \{x \mid x \geqslant 0, x \in \R\}\)，\(N = \{x \mid x^2 < 1, x \in \R\}\)，则 \(M \cap N =\)（\quad）

\choices
    {\([0,1]\)}
    {\((0,1)\)}
    {\((0,1]\)}
    {\([0,1)\)}
\begin{answer}
D
\end{answer}
\end{problem}

\begin{problem}
函数 \( f(x) = \cos \left(2x + \frac{\pi}{4}\right) \) 的最小正周期是（\quad）

\choices
    {\(\frac{\pi}{2}\)}
    {\(\pi\)}
    {\(2\pi\)}
    {\(4\pi\)}
\begin{answer}
B
\end{answer}
\end{problem}

\begin{problem}
已知复数 \(z = 2 - \i\)，则 \(z \cdot \overline{z}\) 的值为（\quad）

\choices
    {\(5\)}
    {\(\sqrt{5}\)}
    {\(3\)}
    {\(\sqrt{3}\)}
\begin{answer}
A
\end{answer}
\end{problem}

\begin{problem}
根据如图所示的框图，对大于 2 的整数 \(N\)，输出的数列的通项公式是（\quad）

\texfigure[max width=0.30\linewidth]{img_repaint/2014/shaanxi_liberal/q4_fig1.tex}

\choices
    {\(a_{n} = 2n\)}
    {\(a_{n} = 2(n - 1)\)}
    {\(a_{n} = 2^{n}\)}
    {\(a_{n} = 2^{n - 1}\)}
\begin{answer}
C
\end{answer}
\end{problem}

\begin{problem}
将边长为 1 的正方形以其一边所在直线为旋转轴旋转一周，所得几何体的侧面积为（\quad）

\choices
    {\(4\pi\)}
    {\(3\pi\)}
    {\(2\pi\)}
    {\(\pi\)}
\begin{answer}
C
\end{answer}
\end{problem}

\begin{problem}
从正方形四个顶点及其中心这5个点中，任取2个点，则这2个点的距离不小于该正方形边长的概率为（\quad）

\choices
    {\(\frac{1}{5}\)}
    {\( \frac{2}{5} \)}
    {\( \frac{3}{5} \)}
    {\(\frac{4}{5}\)}
\begin{answer}
C
\end{answer}
\end{problem}

\begin{problem}
下列函数中，满足“\(f(x + y) = f(x)f(y)\)”的单调递增函数是（\quad）

\choices
    {\(f(x) = x^{\frac{1}{2}}\)}
    {\(f(x) = x^{3}\)}
    {\(f(x) = \left(\frac{1}{2}\right)^x\)}
    {\(f(x) = 3^{x}\)}
\begin{answer}
D
\end{answer}
\end{problem}

\begin{problem}
原命题为“若 \(\frac{a_n + a_{n+1}}{2} < a_n\)，\(n \in \mathbf{N}_+\)，则 \(\{a_n\}\) 为递减数列”，关于逆命题，否命题，逆否命题真假性的判断依次如下，正确的是（\quad）

\choices
    {真，真，真}
    {假，假，真}
    {真，真，假}
    {假，假，假}
\begin{answer}
B
\end{answer}
\end{problem}

\begin{problem}
某公司 10 位员工的月工资（单位：元）为 \(x_{1}\)，\(x_{2}\)，\(\cdots\)，\(x_{10}\)，其均值和方差分别为 \(\overline{x}\) 和 \(s^{2}\)，若从下月起每位员工的月工资增加 100 元，则这 10 位员工下月工资的均值和方差分别为（\quad）

\choices
    {\(\overline{x}\)，\(s^2 + 100^2\)}
    {\(\overline{x} + 100\)，\(s^2 + 100^2\)}
    {\(\overline{x}\)，\(s^2\)}
    {\(\overline{x} + 100\)，\(s^2\)}
\begin{answer}
D
\end{answer}
\end{problem}

\begin{problem}
如图，修建一条公路需要一段环湖弯曲路段与两条直道平滑连接（相切），已知环湖弯曲路段为某三次函数图象的一部分，则该函数的解析式为（\quad）

\bitmapfigure[max width=0.46\linewidth]{img/2014/shaanxi_liberal/q10_fig1.png}

\choices
    {\( y = \frac{1}{2} x^3 - \frac{1}{2} x^2 - x \)}
    {\( y = \frac{1}{2} x^3 + \frac{1}{2} x^2 - 3x \)}
    {\( y = \frac{1}{4} x^3 - x \)}
    {\( y = \frac{1}{4} x^3 + \frac{1}{2} x^2 - 2x \)}
\begin{answer}
A
\end{answer}
\end{problem}

\section{填空题}

\begin{problem}
抛物线 \( y^{2}=4x \) 的准线方程为\fillinblank{}.
\begin{answer}
\(x=-1\)
\end{answer}
\end{problem}

\begin{problem}
已知 \( 4^{a}=2 \)，\( \lg x=a \)，则 \(x=\)\fillinblank{}.
\begin{answer}
\(\sqrt{10}\)
\end{answer}
\end{problem}

\begin{problem}
设 \(0 < \theta < \frac{\pi}{2}\)，向量 \(\bs{a} = (\sin 2\theta, \cos \theta)\)，\(\bs{b} = (1, -\cos \theta)\)，若 \(\bs{a} \cdot \bs{b} = 0\)，则 \(\tan \theta =\)\fillinblank{}.
\begin{answer}
\(\frac{1}{2}\)
\end{answer}
\end{problem}

\begin{problem}
已知 \( f(x) = \frac{x}{1 + x} \)，\( x \geqslant 0 \)，若 \( f_1(x) = f(x) \)，\( f_{n+1}(x) = f(f_n(x)) \)，\( n \in \mathbf{N}_+ \)，则 \( f_{2014}(x) \) 的表达式为\fillinblank{}.
\begin{answer}
\(\frac{x}{2014x+1}\)
\end{answer}
\end{problem}

\section{选考题}

\begin{problem}[A]
设 \(a\)，\(b\)，\(m\)，\(n \in \R\)，且 \(a^2 + b^2 = 5\)，\(ma + nb = 5\)，则 \(\sqrt{m^{2} + n^{2}}\) 的最小值为\fillinblank{}.
\begin{answer}
\(\sqrt{5}\)
\end{answer}
\end{problem}

\reuseproblemnumber
\begin{problem}[B]
如图，\(\triangle ABC\) 中，\(BC = 6\)，以 \(BC\) 为直径的半圆分别交 \(AB\)，\(AC\) 于点 \(E\)，\(F\)，若 \(AC = 2AE\)，则 \(EF =\fillinblank{}.\)

\bitmapfigure[max width=0.42\linewidth]{img/2014/shaanxi_liberal/q15_fig1.png}
\begin{answer}
\(3\)
\end{answer}
\end{problem}

\reuseproblemnumber
\begin{problem}[C]
在极坐标系中，点 \(\left(2, \frac{\pi}{6}\right)\) 到直线 \(\rho \sin \left(\theta - \frac{\pi}{6}\right) = 1\) 的距离是\fillinblank{}.
\begin{answer}
\(1\)
\end{answer}
\end{problem}

\section{解答题}

\begin{problem}
\(\triangle ABC\) 的内角 \(A\)，\(B\)，\(C\) 所对的边分别为 \(a\)，\(b\)，\(c\).

\begin{enumerate}
\item 若 \(a\)，\(b\)，\(c\) 成等差数列，证明：\(\sin A + \sin C = 2\sin (A + C)\)；
\item 若 \(a\)，\(b\)，\(c\) 成等比数列，且 \(c = 2a\)，求 \(\cos B\) 的值.
\end{enumerate}
\begin{solution}
\(\cos B=\frac{3}{4}\).
\end{solution}
\end{problem}

\begin{problem}
四面体 \(ABCD\) 及其三视图如图所示，过棱 \(AB\) 的中点 \(E\) 作平行于 \(AD\)，\(BC\) 的平面分别交四面体的棱 \(BD\)，\(DC\)，\(CA\) 于点 \(F\)，\(G\)，\(H\).
\begin{enumerate}
\item 求四面体 \(ABCD\) 的体积；
\item 证明：四边形 \(EFGH\) 是矩形.
\end{enumerate}

\bitmapfigure[max width=0.55\linewidth]{img/2014/shaanxi_liberal/q17_fig1.png}
\begin{solution}
\begin{enumerate}
\item \(\frac{2}{3}\)；
\item 四边形 \(EFGH\) 是矩形.
\end{enumerate}
\end{solution}
\end{problem}

\begin{problem}
在直角坐标系 \(xOy\) 中，已知点 \(A(1,1)\)，\(B(2,3)\)，\(C(3,2)\). 点 \(P(x,y)\) 在 \(\triangle ABC\) 三边围成的区域（含边界）上，且 \(\overrightarrow{OP} = m\overrightarrow{AB} + n\overrightarrow{AC}\) （\(m\)，\(n \in \R\)）.

\begin{enumerate}
\item 若 \(m = n = \frac{2}{3}\)，求 \(\left|\overrightarrow{OP}\right|\).
\item 用 \(x\)，\(y\) 表示 \(m - n\)，并求 \(m - n\) 的最大值.
\end{enumerate}
\begin{solution}
\begin{enumerate}
\item \(2\sqrt{2}\)；
\item \(m-n=y-x\)，且 \(m-n\) 的最大值为 \(1\).
\end{enumerate}
\end{solution}
\end{problem}

\begin{problem}
某保险公司利用简单随机抽样方法，对投保车辆进行抽样，样本车辆中每辆车的赔付结果统计如下：

\begin{center}
\renewcommand{\arraystretch}{1.2}
\begin{tabular}{c|ccccc}
\hline
赔付金额（元） & 0 & 1000 & 2000 & 3000 & 4000 \\
\hline
车辆数（辆） & 500 & 130 & 100 & 150 & 120 \\
\hline
\end{tabular}
\end{center}

\begin{enumerate}
\item 若每辆车的投保金额均为 2800 元，估计赔付金额大于投保金额的概率；
\item 在样本车辆中，车主是新司机的占 \(10\%\)，在赔付金额为 4000 元的样本车辆中，车主是新司机的占 \(20\%\). 估计在已投保车辆中，新司机获赔金额为 4000 元的概率.
\end{enumerate}
\begin{solution}
\begin{enumerate}
\item \(0.27\)；
\item \(0.24\).
\end{enumerate}
\end{solution}
\end{problem}

\begin{problem}
已知椭圆 \(\frac{x^2}{a^2} + \frac{y^2}{b^2} = 1 \) （\(a > b > 0\)）经过点 \((0, \sqrt{3})\)，离心率为 \(\frac{1}{2}\)，左右焦点分别为 \(F_1(-c, 0)\)，\(F_2(c, 0)\).

\begin{enumerate}
\item 求椭圆的方程.
\item 若直线 \(l: y = -\frac{1}{2} x + m\) 与椭圆交于 \(A\)，\(B\) 两点，与以 \(F_{1}F_{2}\) 为直径的圆交于 \(C\)，\(D\) 两点，且满足 \(\frac{|AB|}{|CD|} = \frac{5\sqrt{3}}{4}\)，求直线 \(l\) 的方程.
\end{enumerate}

\bitmapfigure[max width=0.42\linewidth]{img/2014/shaanxi_liberal/q20_fig1.png}
\begin{solution}
\begin{enumerate}
\item \(\frac{x^2}{4}+\frac{y^2}{3}=1\)；
\item \(y=-\frac{1}{2}x+\frac{\sqrt{3}}{3}\) 或 \(y=-\frac{1}{2}x-\frac{\sqrt{3}}{3}\).
\end{enumerate}
\end{solution}
\end{problem}

\begin{problem}
设函数 \(f(x) = \ln x + \frac{m}{x}\)，\(m \in \R\).

\begin{enumerate}
\item 当 \(m = \e\)（\(\e\) 为自然对数的底数）时，求 \(f(x)\) 的极小值；
\item 讨论函数 \( g(x) = f'(x) - \frac{x}{3} \) 零点的个数；
\item 若对任意 \(b > a > 0\)，\(\frac{f(b) - f(a)}{b - a} < 1\) 恒成立，求 \(m\) 的取值范围.
\end{enumerate}
\begin{solution}
\begin{enumerate}
\item 极小值为 \(2\)；
\item 当 \(m>\frac{2}{3}\) 时有 \(0\) 个零点；当 \(m=\frac{2}{3}\) 或 \(m\leqslant0\) 时有 \(1\) 个零点；当 \(0<m<\frac{2}{3}\) 时有 \(2\) 个零点；
\item \(m\in\left[\frac{1}{4},+\infty\right)\).
\end{enumerate}
\end{solution}
\end{problem}
